\documentclass[12pt,letterpaper]{article}
\usepackage[utf8]{inputenc}
\usepackage[spanish]{babel}
\usepackage{amsmath}
\usepackage{amsfonts}
\usepackage{amssymb}
\usepackage[margin=2.5cm]{geometry}

\begin{document}

\begin{center}
    \Large \textbf{Evaluación de Examen 2: Unidad II} \\
    \large Física para Ciencias de la Salud (005-1713) \\
    \vspace{0.5cm}
    \textbf{Estudiante:} Maria Fernanda Frontado \quad \textbf{C.I.:} 32.747.212
    \vspace{0.3cm} \\
    \large \textbf{Calificación Final: 10/10 puntos}
\end{center}

\vspace{0.5cm}

\section*{Análisis de Resultados}

\subsection*{1. Cinemática (Aceleración media)}
\textbf{Procedimiento y resultado:} Extrae de manera ordenada los datos y plantea correctamente la ecuación de aceleración media ($a = \frac{V_1 - V_0}{t}$). Efectúa la sustitución y división matemática exacta ($0,6 / 0,15 = 4$), logrando el resultado correcto de $4 \text{ m/s}^2$ e indicando muy bien la derivación de sus unidades. \\
\textbf{Veredicto:} \textbf{Correcto} (2/2 pts).

\subsection*{2. Dinámica (Leyes de Newton)}
\textbf{Procedimiento y resultado:} Muestra un planteamiento impecable partiendo de la Segunda Ley de Newton ($F = m \cdot a$) y logrando el despeje algebraico correcto para la aceleración ($a = \frac{F}{m}$). Sustituye y divide $150 \text{ N} / 25 \text{ kg}$, obteniendo exitosamente $6 \text{ m/s}^2$. \\
\textbf{Veredicto:} \textbf{Correcto} (2/2 pts).

\subsection*{3. Trabajo Mecánico}
\textbf{Procedimiento y resultado:} Extrae los datos e identifica la fórmula general del trabajo. Efectúa correctamente la multiplicación directa de la fuerza aplicada y la distancia recorrida ($400 \text{ N} \cdot 0,8 \text{ m}$). El cálculo aritmético da exactamente $320 \text{ Joules (J)}$. \\
\textbf{Veredicto:} \textbf{Correcto} (2/2 pts).

\subsection*{4. Energía Potencial}
\textbf{Procedimiento y resultado:} Extrae rigurosamente los datos del problema (masa, altura y aceleración de gravedad). Plantea el producto de las tres variables según la ecuación de la energía potencial gravitatoria ($E_p = m \cdot g \cdot h$). Obtiene el resultado matemático exacto y verificado de $17,64 \text{ J}$. \\
\textbf{Veredicto:} \textbf{Correcto} (2/2 pts).

\subsection*{5. Potencia Mecánica}
\textbf{Procedimiento y resultado:} Extrae los datos y relaciona correctamente el trabajo efectuado entre el intervalo de tiempo, planteando la fracción respectiva ($P = \frac{w}{t}$). El desarrollo aritmético de $75 \text{ J} / 60 \text{ s}$ arroja la cifra correcta de $1,25 \text{ Watts (W)}$. \\
\textbf{Veredicto:} \textbf{Correcto} (2/2 pts).

\vspace{0.5cm}
\section*{Comentarios Generales y Sugerencias}
¡Felicidades por obtener una calificación perfecta! El examen demuestra un excelente grado de comprensión de las leyes físicas de la Unidad II y una sólida habilidad de cálculo aritmético y algebraico.
\begin{itemize}
    \item La separación constante del problema en bloques de \textit{Datos} por un lado, y el respectivo desarrollo analítico a la derecha, evidencia una mente metódica. Este es un hábito de gran valor para la documentación en ciencias e investigación clínica.
    \item Se dominan a la perfección los despejes y el manejo general de las unidades del Sistema Internacional de Unidades (S.I.).
    \item Como un pequeñísimo apunte formal de cara a futuros documentos técnicos: la convención estricta del Sistema Internacional indica que el símbolo de \textit{Watts} se escribe en mayúscula (\text{W}) por provenir del apellido de un científico (James Watt). Se aplicó excelentemente con los Newtons (N) y los Joules (J).
    \item ¡Sigue con este mismo nivel de dedicación!
\end{itemize}

\end{document}